% BMSBoruvka.tex
% Converted from the original Word draft to LaTeX.
%
% Compile with XeLaTeX (or LuaLaTeX): xelatex BMSBoruvka.tex
% (The document uses Unicode characters directly -- Borůvka's ů, ξ, θ,
%  superscript/subscript glyphs -- so it relies on fontspec + a system
%  Unicode font rather than pdfTeX's 8-bit font model. If you need
%  pdflatex specifically, the Unicode characters below would need to be
%  replaced with LaTeX-native equivalents, e.g. \v{u} or math-mode symbols.)

\documentclass[11pt]{article}

\usepackage{geometry}
\geometry{letterpaper, margin=1in}

\usepackage{amsmath,amssymb}
\usepackage{fontspec}
\usepackage{unicode-math}
\setmainfont{DejaVu Serif}
\setmathfont{DejaVu Math TeX Gyre}

\usepackage{graphicx}
\usepackage{calc}
\usepackage{booktabs,array,longtable}
\usepackage{enumitem}
\usepackage[hidelinks]{hyperref}
\usepackage{xurl}
\urlstyle{same}

% ---- Theorem-like environments (separate counters, matching the
%      plain-number references to "Theorem 2", "Lemma 1", etc. in the prose) ----
\usepackage{amsthm}
\theoremstyle{plain}
\newtheorem{theorem}{Theorem}
\newtheorem{lemma}{Lemma}
\theoremstyle{remark}
\newtheorem{remark}{Remark}

\setlength{\emergencystretch}{3em}
\providecommand{\tightlist}{\setlength{\itemsep}{0pt}\setlength{\parskip}{0pt}}

\title{BMSBorůvka: A Multi-Level Bucket Framework for Minimum Spanning Tree}
\author{}
\date{}

\begin{document}

\maketitle







\begin{abstract}
The minimum spanning tree (MST) problem traditionally relies on global
edge sorting, leading to an \(O(m\ log\ n)\) time complexity barrier in
algorithms such as Kruskal\textquotesingle s Algorithm and
Prim\textquotesingle s Algorithm. In this paper, we propose BMSBorůvka,
a bucket-based Borůvka-style algorithm that avoids full sorting through
phased edge activation and multi-level graph compression. Our approach
achieves \(O(m\ \log^{\frac{2}{3}}\ n)\) time on sparse graphs while
maintaining \(O(m\ log\ n)\) performance in the general case.
\end{abstract}

\section{Introduction}

A minimum spanning tree (MST) or minimum weight spanning tree is a
subset of the edges of a connected, edge-weighted undirected graph that
connects all the vertices together, without any cycles and with the
minimum possible total edge weight. Classical algorithms such as Prim's
and Kruskal's algorithms achieve optimal or near-optimal performance in
general settings, but they rely heavily on comparison-based
mechanisms---either through repeated priority queue operations in Prim's
algorithm or through explicitly sorting all edges by weight in Kruskal's
algorithm. As a result, their running time is dominated by
comparison-based sorting costs, typically on the order of
\(O(m\ log\ n)\), where n is the number of vertices in the graph.

This reliance reflects a broader sorting bottleneck that also appears in
shortest path algorithms. It has long been known that when values can be
classified in constant time, sorting can be performed in linear time,
avoiding comparison-based ordering. Recent work {[}DMMS+25{]} on
breaking the sorting barrier for directed single-source shortest paths
exploits this idea by replacing comparison-heavy priority queue
operations with classification-based techniques using multi-level
buckets and graph sparsification. Motivated by this approach, we extend
similar ideas to the Borůvka framework and demonstrate improved
performance.

\section{Related Work}

The minimum spanning tree (MST) problem is one of the most fundamental
problems in graph algorithms, with a long line of research focusing on
improving its time complexity. Classical algorithms such as Kruskal
{[}Kru56{]} and Prim~{[}Pri57{]} algorithms achieve \(O(m\ log\ n)\)
time by relying on comparison-based sorting or priority queue
operations.

A major line of work has sought to reduce or bypass the sorting
bottleneck inherent in these approaches. Chazelle {[}Cha00{]} gave a
deterministic MST algorithm running in \(O(m\ \alpha(m,\ n))\) time,
where α is the inverse Ackermann function. His algorithm avoids explicit
global sorting by using soft heaps, effectively allowing controlled key
corruption to achieve near-linear performance.

Randomization has also played a crucial role in overcoming the sorting
barrier. Karger, Klein, and Tarjan {[}KKT95{]} presented a randomized
linear-time MST algorithm with expected running time \(O(m)\), combining
random sampling, filtering, and contraction techniques. This result
demonstrates that, in the randomized setting, the sorting barrier can be
circumvented entirely.

Pettie and Ramachandran {[}PR02{]} further advanced the theoretical
understanding of MST by designing an optimal comparison-based algorithm.
Their algorithm achieves a running time of \(O(m\ \alpha(m,\ n))\) plus
a term matching the decision-tree complexity of MST, thereby showing
that MST can be solved optimally in the comparison model without
requiring explicit sorting.

Despite these advances, most near-linear or optimal algorithms rely on
sophisticated data structures, randomization, or implicit sorting
mechanisms. In contrast, recent work on single-source shortest paths
(SSSP) has explored breaking the sorting barrier through hierarchical
bucketing and sparsification techniques {[}DMMS+25{]}. Inspired by these
developments, our work investigates whether similar ideas can be applied
to MST construction.

In particular, we propose a bucket-based multi-level sparsification
framework that groups edges by weight ranges and progressively reduces
the graph while preserving candidate MST edges. This approach avoids
global edge sorting and instead relies on structured classification and
local processing, drawing conceptual parallels to Borůvka-style
algorithms. Our goal is to achieve a theoretically sound and practically
efficient alternative that reshapes the role of sorting in MST
algorithms.

\section{Preliminaries}

\subsection{Graphs and the Minimum Spanning Tree Problem}

An undirected weighted graph is a triple \(G = (V,E,w)\) where V is a
finite set of vertices,
\(E\  \subseteq \text{\{\{}u,v\text{\}}\  \mid u,v\  \in V,\ u\  \neq v\text{\}}\ \)is
a set of edges, and \(w:E \rightarrow \mathbb{Z}\) assigns an integer
weight to each edge. We write \(n\  = \ |V|\) and \(m\  = \ |E|\).

A path between two vertices u and v is a sequence of distinct vertices u
= v₀, v₁, \ldots, vₖ = v such that
\(\left\{ v_{i},\ v_{i + 1} \right\} \in \ E\) for each i. A graph is
connected if a path exists between every pair of vertices; otherwise it
is disconnected.

A spanning tree connects all vertices of a graph with no cycles; when
edges carry weights, the spanning tree that minimises the sum of edge
weights is called the minimum spanning tree (MST) {[}Hil24{]}. A
spanning tree of a connected graph G is a subgraph \(T\  \subseteq E\)
that (i) connects all n vertices, (ii) contains no cycles, and (iii) has
exactly \(n - 1\) edges. A minimum spanning tree (MST) is a spanning
tree whose total edge weight \(\sum_{e \in T}^{}{w(e)}\) is minimised
over all spanning trees. If G is disconnected with k connected
components, the analogous structure is a minimum spanning forest (MSF),
which has \(n - k\) edges.

Two classical properties are central to all MST algorithms:

\textbf{Cut Property.} Let \(S\  \subset V\) be any non-empty proper
subset of vertices. The cut defined by S consists of all edges with one
endpoint in S and the other in \(V\  \smallsetminus S\). If e is the
unique minimum-weight edge of this cut, then e belongs to every MST.
More generally, any minimum-weight edge across any cut can be safely
added to a partial MST without creating a suboptimal solution.

\textbf{Cycle Property.} If e is the unique maximum-weight edge on some
cycle in G, then e does not belong to any MST.

These two properties together characterise exactly the edges that belong
to the MST.

\subsection{Borůvka\textquotesingle s Algorithm}

Borůvka\textquotesingle s algorithm {[}Bor26{]} is the foundation on
which BMSBoruvka is built. The key idea is to grow many components of
the MST simultaneously rather than one edge at a time.

Starting with n isolated components (one per vertex), the algorithm
repeats the following Borůvka round until only one component remains:

\textbf{Find cheapest outgoing edge.} For each current component,
identify its minimum-weight edge that crosses to a different component
(its cheapest outgoing edge, or COE). Break ties deterministically ---
here by edge index.

\textbf{Merge.} Add every selected COE to the MST and merge the two
components it connects. Because multiple components may select the same
edge, a Disjoint Set Union structure (Section 3.4) tracks which
components have already merged.

After one round, every component that had an outgoing edge has merged
with at least one neighbour, so the number of components at least halves
per round. Starting from n components and halving each round, the
algorithm terminates in at most ⌈log₂ n⌉ rounds. Each round scans all m
edges, giving a total time of \(O(m\ log\ n)\) for a connected graph.

\subsection{Graph Compression (Super-nodes and Super-edges)}

After several Borůvka rounds the number of components has shrunk
significantly --- for example, from 1,000,000 vertices down to a few
hundred "super-nodes". It is wasteful to keep scanning all m original
edges when most of them now connect vertices within the same super-node
and can never be in the MST.

Graph compression solves this by replacing the current graph with a
smaller compressed graph whose vertices are the surviving super-nodes.
More precisely, if two super-nodes A and B are connected by several
original edges, we keep only the minimum-weight one (by the cycle
property, all others are dominated). Edges within a single super-node
become self-loops and are discarded.

After compression the new graph has C vertices (one per remaining
super-node) and at most \(C - 1\) edges relevant to the MST. For sparse
graphs, C shrinks rapidly (see Lemma 1), so subsequent Borůvka rounds
operate on a much smaller edge set. This is the primary reason
BMSBoruvka outperforms plain Borůvka on sparse graphs.

The algorithm maintains a mapping origToComp{[}v{]} that records, for
each original vertex v, the ID of the super-node it currently belongs
to. This allows new edges to be pushed into the compressed graph without
re-reading the original adjacency list.

\subsection{Disjoint Set Union (DSU)}

A Disjoint Set Union (DSU), also called Union-Find, is a data structure
that efficiently tracks which elements belong to the same set
(component). It supports two operations:

\textbf{find(x)} --- returns the representative (root) of the set
containing x. Two elements are in the same component if and only if
their representatives are equal.

\textbf{unite(a, b)} --- merges the sets containing a and b, returning
the new representative.

This implementation uses two standard optimisations. Union by rank
{[}Tar75{]} always attaches the shorter tree under the taller one,
keeping trees shallow. Path halving {[}TvL84{]} during find() rewires
every other node on the path directly to its grandparent, compressing
future traversals. Together these give an amortised cost of
\(O(\alpha(n))\) per operation, where α is the inverse Ackermann
function --- a value no larger than 4 for any input that fits in
computer memory. For all practical purposes this is \(O(1)\).

Crucially, find() is not a comparison in the graph-theory sense --- it
performs pointer traversals through the parent{[}{]} array. This
distinction matters for the cost model (Section 3.6).

\subsection{Selection vs. Sorting: nth\_element}

A classical result in algorithm theory is that selecting the k-th
smallest element of an array of n items requires only \(O(n)\)
comparisons in the worst case, whereas sorting the entire array requires
\(\Omega\left( n\log n \right)\) comparisons. This gap of a factor log n
is exactly what BMSBoruvka exploits.

The C++ standard library function std::nth\_element(first, nth, last,
cmp) rearranges the range {[}first, last) so that:

\begin{itemize}
\item
  The element at position nth is the element that would be there in
  sorted order.
\item
  Every element before nth is no greater than the element at nth.
\item
  Every element after nth is no less than the element at nth.
\item
  Elements before nth and elements after nth are not sorted relative to
  each other.
\end{itemize}

The standard library guarantees worst-case \(O(n)\) time for this
operation using an algorithm called Introselect, which combines the
average-case-fast Quickselect with a worst-case-optimal
Median-of-Medians fallback. Crucially, the bound is unconditional --- it
holds regardless of the distribution or repetition of edge weights.

In BMSBoruvka, nth\_element is called once per super-step on the
unactivated suffix of the edge array. It moves the lightest k edges to
the front (unordered) so they can be processed in that super-step, while
deferring heavier edges to later super-steps. This preserves the
weight-prefix property needed for correctness (Section 4.2) without
paying the \(O(m\ log\ m)\) cost of a full global sort.

\subsection{Cost Model}

We analyse BMSBoruvka under the unit-cost word RAM model. In this model
every elementary operation on a machine word --- integer arithmetic,
array indexing, pointer dereference, hash-table lookup --- costs exactly
1 unit of time, regardless of its type. All integers, edge weights, and
array indices fit in a single word.

Under this model:

\begin{itemize}
\item
  An edge-weight comparison costs 1 unit.
\item
  A DSU find() call costs \(O(\alpha(n))\) units (amortised over any
  sequence of operations), which we treat as \(O(1)\) throughout the
  analysis since \(\alpha(n) \leq 4\) in practice.
\item
  A hash\_map insertion or lookup costs \(O(1)\) units amortised.
\item
  Allocating or writing a block of k memory words costs \(O(k)\) units.
\end{itemize}

This model is more general than the pure comparison model used in
classical MST lower bounds (which counts only edge-weight comparisons).
The distinction matters here because Phase B spends roughly four times
as many units on DSU operations as on comparisons, and Phase C
additionally incurs hash-table costs. The RAM model captures all of
these.

The RAM model does not capture cache effects. In practice, find()
accesses the parent{[}{]} array in a graph-dependent, near-random order;
for large n (e.g., n = 10\^{}6, parent{[}{]} occupying 4 MB) this causes
L2/L3 cache pressure that is invisible to the RAM model. This explains
why the empirical throughput metric ̂t (operations per millisecond)
trends upward from n = 10,000 to n = 100,000 before stabilising: at
small n the DSU tree is shallower and the working set fits in cache; at
large n both effects plateau together. The asymptotic bounds derived
under the RAM model are confirmed by the empirical stabilisation of
\(\widehat{t}\) for \(n\  \geq \ 100,000\).

\subsection{Summary of Notation}

The following symbols are used consistently throughout Sections 4 and 5
as shown in Table 1.

Table 1. Summary of Notation.

\begin{longtable}[]{@{}
  >{\raggedright\arraybackslash}p{(\columnwidth - 2\tabcolsep) * \real{0.2575}}
  >{\raggedright\arraybackslash}p{(\columnwidth - 2\tabcolsep) * \real{0.7425}}@{}}
\toprule\noalign{}
\begin{minipage}[b]{\linewidth}\raggedright
Symbol
\end{minipage} & \begin{minipage}[b]{\linewidth}\raggedright
Meaning
\end{minipage} \\
\midrule\noalign{}
\endhead
\bottomrule\noalign{}
\endlastfoot
\(G\  = \ (V,\ E,\ w)\) & Input graph: vertex set \(V\), edge set \(E\),
weight function \(w:E \rightarrow \mathbb{Z}\) \\
\(n\) & Number of vertices, \(n\  = \ |V|\) \\
\(m\) & Number of edges, \(m\  = \ |E|\) \\
MST / MSF & Minimum spanning tree (connected graph) or minimum spanning
forest (disconnected) \\
\(t\) & Number of Borůvka rounds per super-step:
\(t\  = \ \left\lceil \log_{2}^{\frac{2}{3}}(n) \right\rceil\ \) \\
\(P\) & Number of super-steps:
\(P = \left\lceil \log_{2}(n)\text{/}t \right\rceil = \Theta\left( \log^{1\text{/}3}n \right)\) \\
\(k\) & Edges activated per super-step:
\(k = \left\lfloor m/P \right\rfloor\) \\
\(C_{i}\) & Number of super-nodes (components) at the start of
super-step i; \(C_{0} = n\) \\
\(M_{i}\) & Size of curEdges entering Phase B of super-step i \\
\(m_{i}^{'}\) & Number of distinct inter-component edges after
compress() in super-step i \\
origToComp{[}v{]} & Maps original vertex v to its current super-node
ID \\
curEdges & Working edge list in the compressed graph, updated each
super-step \\
\(\alpha(n)\) & Inverse Ackermann function; \(\alpha(n)\  \leq \ 4\) for
all practical n \\
\(O( \cdot )\) & Standard asymptotic upper bound (worst-case unless
stated otherwise) \\
\end{longtable}

\section{Results}

\subsection{The Algorithm}

The BMSBoruvka algorithm operates on an undirected weighted graph
\(G = (V,\ E)\) with \(n\  = \ |V|\) vertices and \(m\  = \ |E|\) edges.
Two integer parameters govern the structure of the computation:

\begin{quote}
\(t\  = \ \left\lceil \log_{2}^{\frac{2}{3}}(n) \right\rceil\  \geq \ 1\)
(Borůvka rounds per super-step)

\(P = \left\lceil \log_{2}(n)\text{/}t \right\rceil = \Theta\left( \log^{1\text{/}3}n \right)\)
(number of super-steps)
\end{quote}

The algorithm proceeds in P super-steps. Each super-step consists of
three phases:

\begin{enumerate}
\def\labelenumi{\arabic{enumi}.}
\item
  \textbf{Phase A (Activation).} A call to nth\_element partitions the
  as-yet unactivated suffix of edges{[}{]} so that the k = ⌊m/P⌋
  lightest edges move to the front (unordered among themselves) and
  every remaining edge is no lighter than any activated edge. These are
  pushed into curEdges via the origToComp{[}{]} mapping. Because the
  activated prefix is always a weight-prefix of E, the cut property is
  preserved.
\item
  \textbf{Phase B (Borůvka Rounds).} Execute t standard Borůvka rounds
  on curEdges. In each round every component selects its minimum-weight
  outgoing edge, all selected edges are added to the MST, and edges
  whose endpoints have been merged are filtered out.
\item
  \textbf{Phase C (Compression).} compress() renumbers the surviving
  super-nodes 0, \ldots, Cᵢ₊₁−1, rebuilds curEdges with the new IDs
  (discarding self-loops), and updates origToComp{[}0..n). The DSU is
  then reinitialised in the compressed space.
\end{enumerate}

The outer loop terminates when all n−1 MST edges have been collected,
when uf.comps drops to 1, or after at most P super-steps.

\subsection{Correctness Analysis}

\begin{theorem}[MST / MSF Correctness]
BMSBoruvka returns a
minimum spanning tree of G if G is connected, and a minimum spanning
forest otherwise.
\end{theorem}

\begin{proof}
We verify the three invariants that together guarantee
correctness.
\end{proof}

\textbf{Cut property preservation.} nth\_element guarantees that
edges{[}curActivated..activateEnd) contains the globally lightest k
unactivated edges: every edge in the window has weight ≤ every edge
outside it. Because curActivated advances monotonically across
super-steps, the cumulative activated set is always a weight-prefix of
E. The cut property requires only that the minimum-weight crossing edge
of each cut be present in curEdges at the moment that cut is resolved;
since the activated set is a weight-prefix, this condition is always
satisfied.

\textbf{Phase B safety.} Each Borůvka round selects, for every active
component, its cheapest outgoing edge. By the cut property of matroids,
each such edge is safe (i.e., belongs to some MST). No edge is selected
twice because self-loops are filtered before each round.

\textbf{Compression correctness.} Every edge in the compressed graph
corresponds to a unique original edge via orig\_eidx. Self-loops (edges
whose endpoints lie in the same component) are discarded, which is sound
because a self-loop cannot belong to any spanning tree. The
origToComp{[}{]} mapping is updated consistently so that future Phase A
activations route edges to the correct super-nodes.

For disconnected graphs the outer loop terminates once no component has
an outgoing edge in curEdges; the result is a minimum spanning forest
with \(n - k\) edges, where k is the number of connected components.
Correctness was independently verified by 3,000 random stress tests
against Kruskal's algorithm.

\subsection{Time Complexity Analysis}

\begin{lemma}[Component Shrinkage]
After super-step i, the
component count satisfies \(C_{i + 1} \leq C_{i}/2^{t}\). In particular,
\(C_{P}\  = \ 1\) for a connected graph, and
\(C_{i} \leq \frac{n}{2^{it}}\).
\end{lemma}

\begin{proof}
Each of the t Borůvka rounds in Phase B selects, for every
component, its minimum-weight outgoing edge; by the cut property every
selected edge is safe. After one such round every component that
possesses an outgoing edge merges with at least one neighbour, so the
component count at least halves. Repeating t times gives
\(C_{i + 1} \leq \frac{C_{i}}{2^{t}}\). Since
\(t \cdot P = \Theta\left( \log n \right)\), after P super-steps we have
\(C_{P} \leq \frac{n}{2^{tP}} \rightarrow 1\).
\end{proof}

\begin{lemma}[nth\_element Total Cost]
The aggregate cost of all
nth\_element calls is \(O(mP) = O\left( m\log^{1\text{/}3}n \right)\).
\end{lemma}

\begin{proof}
Super-step i operates on the suffix of size
\(m\  - \ i \cdot k\); introselect guarantees worst-case \(O(size)\)
time independent of weight distribution. Summing the arithmetic
progression over steps 0 through P−2 (the final step activates all
remaining edges without a partition call):
\(\sum_{i = 0}^{P - 2}(m - i \cdot k)\  \leq \ \frac{m(P + 1)}{2} = O(mP) = O\left( m\log^{1\text{/}3}n \right)\).
No degenerate case exists; the bound is unconditional.
\end{proof}

\begin{lemma}[Phase A Total Cost]
The aggregate cost of all Phase A
activation loops is \(O(m)\).
\end{lemma}

\begin{proof}
Each original edge is pushed into curEdges at most once
(curActivated advances monotonically) and each origToComp lookup is
\(O(1)\). The total work across all super-steps is therefore \(O(m)\).
\end{proof}

\begin{lemma}[Phase B, General Graphs]
The total cost of all Phase
B Borůvka rounds is \(O(m\ log\ n)\).
\end{lemma}

\begin{proof}
Let \(M_{i}^{(r)}\) denote the size of curEdges at the
start of round r within super-step i. Each round scans curEdges in
\(O(M_{i}^{(r)})\) time and performs DSU operations in \(O(C_{i})\).
Since each edge survives at most t rounds before its endpoints merge,
\(\sum_{r}^{}M_{i}^{(r)} \leq t \cdot M_{i}\). Applying the crude bound
\(M_{i} \leq \ m\) and summing over all super-steps:
\(T_{\text{Phase B}} \leq O(t \cdot P \cdot m) = O\left( \log^{2\text{/}3}(n) \cdot \log^{1\text{/}3}(n) \cdot m \right) = O\left( m\log n \right)\).
\end{proof}

\begin{lemma}[Phase B, Sparse Graphs]
For sparse graphs with
\(m = O(n)\), the total Phase B cost is
\(O\left( n\log^{2\text{/}3}n \right)\).
\end{lemma}

\begin{proof}
Under the sparse assumption
\(m_{i}^{'} = O\left( C_{i} \right)\) (the compressed graph retains
\(O(1)\) average super-node degree), we have
\(M_{i} \leq m_{i}^{'} + k = O\left( C_{i} + k \right)\). Summing:
\(\sum_{i}^{}M_{i} \leq \sum_{i}^{}C_{i} + m\). By Lemma 1,
\(\sum_{i}^{}C_{i} \leq \frac{n}{1 - 2^{- t}} \leq 2n\). With
\(m\  = \ O(n)\), \(\sum_{i}^{}M_{i} = O(n)\), so the total Phase B cost
is \(O(t \cdot n) = O\left( n\log^{2\text{/}3}n \right)\).
\end{proof}

\begin{lemma}[Phase C Total Cost]
The total cost of all compress()
calls is \(O\left( \sum_{i}^{}M_{i} + P \cdot n \right)\).
\end{lemma}

\begin{proof}
Each compress() call performs: (i) an \(O(M_{i})\) scan to
filter self-loops; (ii) \(O(C_{i})\) work to enumerate roots and build
the renaming map; (iii) \(O\left( m_{i}^{'} \right)\) edge renaming;
(iv) an \(O(n)\) origToComp update; (v) an \(O(C_{i} + C_{i + 1})\)
array reset. Summing over P super-steps yields
\(O\left( \sum_{i}^{}M_{i} + P \cdot n \right)\). For general graphs
this is \(O\left( m\log^{1\text{/}3}n \right)\); for sparse graphs,
\(O\left( n\log^{1\text{/}3}n \right)\).
\end{proof}

\begin{theorem}[Sparse Graphs, m = O(n)]
BMSBoruvka computes
an MST in time \(T = O\left( n\log^{2\text{/}3}n \right)\).
\end{theorem}

\begin{theorem}[General Graphs]
BMSBoruvka computes an MST
in time \(T = O(m\ log\ n)\).
\end{theorem}

\begin{proof}[Proof of Theorem 2]
Collecting costs for \(m\  = \ O(n)\):
\end{proof}

Table 2. Cost breakdown of BMSBoruvka on sparse graphs.

\begin{longtable}[]{@{}
  >{\raggedright\arraybackslash}p{(\columnwidth - 4\tabcolsep) * \real{0.3846}}
  >{\raggedright\arraybackslash}p{(\columnwidth - 4\tabcolsep) * \real{0.3419}}
  >{\raggedright\arraybackslash}p{(\columnwidth - 4\tabcolsep) * \real{0.2735}}@{}}
\toprule\noalign{}
\begin{minipage}[b]{\linewidth}\raggedright
Component
\end{minipage} & \begin{minipage}[b]{\linewidth}\raggedright
Cost (sparse, m = O(n))
\end{minipage} & \begin{minipage}[b]{\linewidth}\raggedright
Dominant?
\end{minipage} \\
\midrule\noalign{}
\endhead
\bottomrule\noalign{}
\endlastfoot
nth\_element (Lemma 2) & \(O\left( n\log^{1\text{/}3}n \right)\) & No \\
Phase A activation (Lemma 3) & \(O(n)\) & No \\
Phase B Borůvka (Lemma 5) & \(O\left( n\log^{2\text{/}3}n \right)\) &
Yes \\
Phase C compress (Lemma 6) & \(O\left( n\log^{1\text{/}3}n \right)\) &
No \\
Total & \(O\left( n\log^{2\text{/}3}n \right)\) & \\
\end{longtable}

\begin{proof}
Phase B dominates since
\(\log^{1\text{/}3}(n) \leq \log^{2\text{/}3}(n)\) for all \(n \geq 2\),
so every other component lies in
\(O\left( n\log^{1\text{/}3}n \right) \subseteq O\left( n\log^{2\text{/}3}n \right)\).
\end{proof}

\begin{proof}[Proof of Theorem 3]
For general graphs, apply the crude bound
\(M_{i}\  \leq \ m\) in Lemma 4:
\end{proof}

Table 3. Cost breakdown of BMSBoruvka on general graphs.

\begin{longtable}[]{@{}
  >{\raggedright\arraybackslash}p{(\columnwidth - 4\tabcolsep) * \real{0.3846}}
  >{\raggedright\arraybackslash}p{(\columnwidth - 4\tabcolsep) * \real{0.3419}}
  >{\raggedright\arraybackslash}p{(\columnwidth - 4\tabcolsep) * \real{0.2735}}@{}}
\toprule\noalign{}
\begin{minipage}[b]{\linewidth}\raggedright
Component
\end{minipage} & \begin{minipage}[b]{\linewidth}\raggedright
Cost (general)
\end{minipage} & \begin{minipage}[b]{\linewidth}\raggedright
Dominant?
\end{minipage} \\
\midrule\noalign{}
\endhead
\bottomrule\noalign{}
\endlastfoot
nth\_element (Lemma 2) & \(O\left( m\log^{1\text{/}3}n \right)\) & No \\
Phase A activation (Lemma 3) & \(O(m)\) & No \\
Phase B Borůvka (Lemma 4) & \(O(m\ log\ n)\) & Yes \\
Phase C compress (Lemma 6) & \(O\left( m\log^{1\text{/}3}n \right)\) &
No \\
Total & \(O(m\ log\ n)\) & \\
\end{longtable}

\begin{proof}
Phase B dominates:
\(O(t \cdot P \cdot m) = O\left( \log(n) \cdot m \right)\).
\end{proof}

\begin{remark}[Sparse-graph condition]
The bound
\(O\left( n\log^{2\text{/}3}n \right)\) requires \(m\  = \ O(n)\)
(bounded average degree). This is needed to ensure
\(m_{i}^{'} = O\left( C_{i} \right)\) after each compression, which in
turn gives \(\sum_{i}^{}M_{i} = O(n)\). For denser graphs
\(m = O\left( n\, f(n) \right),\ \ f(n) \rightarrow \infty\), Phase B
costs \(O\left( n\log^{2\text{/}3}(n) \cdot f(n) \right)\), which still
improves on the plain Borůvka bound
\(O\left( m\log n \right) = O\left( n\, f(n)\log n \right)\) by a factor
of \(\log^{1\text{/}3}n\).
\end{remark}

\section{Benchmark}

\subsection{Set up}

Binaries were compiled with g++ -O2 -std=c++17, and the parallel variant
additionally with -fopenmp. Three algorithms were evaluated:

\begin{enumerate}
\def\labelenumi{\arabic{enumi}.}
\item
  \textbf{boruvka.} Standard sequential Borůvka, implemented with the
  same DSU and edge-comparison routines as BMSBoruvka. Serves as the
  baseline.
\item
  \textbf{par\_Tk.} OpenMP-parallelised Borůvka with k threads (k ∈ \{2,
  4\}). Phase 1 parallelises the COE scan with thread-local best-edge
  arrays; Phase 2 reduces them sequentially; Phase 3 unions
  sequentially.
\item
  \textbf{bms.} BMSBoruvka (v6, nth\_element preprocessing), the
  algorithm analysed in Section 4.
\end{enumerate}

Test graphs were generated with a custom random graph generator using
seed 42. Every graph was guaranteed to be connected (a random spanning
path was embedded first). Two density classes were evaluated:

\begin{enumerate}
\def\labelenumi{\arabic{enumi}.}
\item
  \textbf{Sparse (m≈5n).}
  \(\mathbf{n \in \{ 10000,\ 100000,\ 500000,1000000\},m = 5}\mathbf{n.}\)
\item
  \textbf{Dense (m≈50n).}
  \(\mathbf{n \in \{ 50000,100000\},m = 50}\mathbf{n.}\)
\end{enumerate}

Timing used time.perf\_counter() in Python with a one warm-up run
discarded before each (algorithm, test) combination. The number of timed
trials was determined adaptively: a probe run estimated the single-run
latency, and trials were set to clamp(⌊24000 ms ÷ budget⌋, 3, 60),
targeting roughly two minutes of total measurement time per test case
across all algorithms. Program output was redirected to /dev/null to
exclude I/O from timings.

\subsection{Set up}

\textbf{Wall-clock time}. Mean, minimum, maximum, and standard deviation
(σ) of timed runs, reported in milliseconds.

\textbf{Shrinkage factor ξ}. To quantify how quickly each algorithm
reduces the graph, we define \(\xi_{i} = \frac{C_{i}}{C_{i + 1}}\),
where \(C_{i}\) is the number of connected components before step i. For
sequential Borůvka, the step is a single round; for BMSBoruvka, the step
is a full super-step. We report \(\xi_{\mu}\), the geometric mean of all
per-step ξ values over a complete run. A larger \(\xi_{\mu}\) indicates
faster component consolidation per step. The theoretical lower bound for
one Borůvka round is \(\xi \geq 2\) (components at least halve); for a
BMSBoruvka super-step the theoretical value is \(\xi = 2^{t}\).

\textbf{Normalised throughput} \(\widehat{\mathbf{t}}\)\textbf{.} To
assess whether the empirical performance matches the theoretical
complexity, we define
\(\widehat{t} = \frac{\theta_{0}(n,m)}{T_{\text{mean}}}\), where
\(T_{mean}\) is the mean wall-clock time in milliseconds and
\(\theta_{0}\) is the algorithm's theoretical operation count:

boruvka / par:
\(\theta_{0} = m \cdot \, log_{2}(n)\quad\left\lbrack O(m\, logn) \right\rbrack\)

bms:
\(\theta_{0} = n \cdot log_{2}^{\frac{2}{3}}(n)\quad\left\lbrack O\left( nlog^{\frac{2}{3}}\, n \right) \right\rbrack\)

If \(\widehat{t}\) converges to a constant as n grows, the theoretical
complexity bound is tight in practice. A rising \(\widehat{t}\)
indicates improving hardware utilisation; a falling \(\widehat{t}\)
would indicate growing overhead relative to the theoretical model.

\subsection{Shrinkage Factor ξ}

Table 4 reports \(\xi_{\mu}\) (averaged over three independent runs) for
the sparse graph series.

Table 4. Measured shrinkage factor \(\xi_{\mu}\) (geometric mean) for
sparse graphs (m ≈ 5n).

\begin{longtable}[]{@{}
  >{\raggedright\arraybackslash}p{(\columnwidth - 12\tabcolsep) * \real{0.1389}}
  >{\raggedright\arraybackslash}p{(\columnwidth - 12\tabcolsep) * \real{0.1496}}
  >{\raggedright\arraybackslash}p{(\columnwidth - 12\tabcolsep) * \real{0.1346}}
  >{\raggedright\arraybackslash}p{(\columnwidth - 12\tabcolsep) * \real{0.1827}}
  >{\raggedright\arraybackslash}p{(\columnwidth - 12\tabcolsep) * \real{0.1165}}
  >{\raggedright\arraybackslash}p{(\columnwidth - 12\tabcolsep) * \real{0.1816}}
  >{\raggedright\arraybackslash}p{(\columnwidth - 12\tabcolsep) * \real{0.0962}}@{}}
\toprule\noalign{}
\begin{minipage}[b]{\linewidth}\raggedright
n
\end{minipage} & \begin{minipage}[b]{\linewidth}\raggedright
m
\end{minipage} & \begin{minipage}[b]{\linewidth}\raggedright
Algorithm
\end{minipage} & \begin{minipage}[b]{\linewidth}\raggedright
\(\xi_{\mu}\) (measured)
\end{minipage} & \begin{minipage}[b]{\linewidth}\raggedright
\[\xi_{theory}\]
\end{minipage} & \begin{minipage}[b]{\linewidth}\raggedright
Steps / Rounds
\end{minipage} & \begin{minipage}[b]{\linewidth}\raggedright
t
\end{minipage} \\
\midrule\noalign{}
\endhead
\bottomrule\noalign{}
\endlastfoot
10,000 & 50,000 & bms & 40.82 & 64 & 2 & 6 \\
10,000 & 50,000 & boruvka & 5.07 & 2 & 5 & --- \\
100,000 & 500,000 & bms & 45.64 & 128 & 2 & 7 \\
100,000 & 500,000 & boruvka & 5.67 & 2 & 6 & --- \\
500,000 & 2,500,000 & bms & 44.19 & 256 & 2 & 8 \\
500,000 & 2,500,000 & boruvka & 5.35 & 2 & 7 & --- \\
1,000,000 & 5,000,000 & bms & 42.84 & 256 & 2 & 8 \\
1,000,000 & 5,000,000 & boruvka & 5.90 & 2 & 7 & --- \\
\end{longtable}

Several observations follow from Table 4. First, BMSBoruvka consistently
terminates in only two super-steps for all sizes tested, compared to
five to seven rounds for sequential Borůvka. This is a direct
consequence of the large per-super-step shrinkage: with t = 8 and
\(\xi_{\mu} \approx 43\) for \(n = 10^{6}\), two super-steps reduce
\(10^{6}\) components to approximately
\(10^{6}\text{/}43^{2} \approx 540\), and the residual is absorbed by
the second step.

Second, the measured \(\xi_{\mu}\) for BMSBoruvka (41--46) is
substantially smaller than the theoretical value \(2^{t}\)(64--256).
This is expected: the theoretical value assumes every super-step runs
all t Borůvka rounds to completion, but in practice the algorithm
terminates early within a super-step as soon as no merges are possible.
The measured \(\xi_{\mu}\) therefore reflects the actual convergence
rate rather than the worst-case guarantee.

Third, sequential Borůvka achieves \(\xi_{\mu} \approx 5 - 6\) per round
on these random sparse graphs, well above the theoretical lower bound of
2. This is characteristic of random graphs, where many components
simultaneously find outgoing edges to previously isolated vertices,
producing merges that exceed the minimum halving guarantee.

\subsection{Wall-Clock Time and Speedup}

Tables 5 and 6 report the mean runtime, standard deviation, and speedup
ratios for sparse and dense graph instances respectively.

Table 5. Runtime results for sparse graphs (m ≈ 5n). `vs boruvka' =
boruvka mean / algo mean; `vs bms' = bms mean / algo mean. Ratios
\textgreater{} 1 indicate the algorithm is faster than the reference.

\begin{longtable}[]{@{}
  >{\raggedright\arraybackslash}p{(\columnwidth - 14\tabcolsep) * \real{0.1395}}
  >{\raggedright\arraybackslash}p{(\columnwidth - 14\tabcolsep) * \real{0.1419}}
  >{\raggedright\arraybackslash}p{(\columnwidth - 14\tabcolsep) * \real{0.0841}}
  >{\raggedright\arraybackslash}p{(\columnwidth - 14\tabcolsep) * \real{0.1398}}
  >{\raggedright\arraybackslash}p{(\columnwidth - 14\tabcolsep) * \real{0.1449}}
  >{\raggedright\arraybackslash}p{(\columnwidth - 14\tabcolsep) * \real{0.0869}}
  >{\raggedright\arraybackslash}p{(\columnwidth - 14\tabcolsep) * \real{0.1571}}
  >{\raggedright\arraybackslash}p{(\columnwidth - 14\tabcolsep) * \real{0.1057}}@{}}
\toprule\noalign{}
\begin{minipage}[b]{\linewidth}\raggedright
n
\end{minipage} & \begin{minipage}[b]{\linewidth}\raggedright
m
\end{minipage} & \begin{minipage}[b]{\linewidth}\raggedright
trials
\end{minipage} & \begin{minipage}[b]{\linewidth}\raggedright
Algorithm
\end{minipage} & \begin{minipage}[b]{\linewidth}\raggedright
Mean (ms)
\end{minipage} & \begin{minipage}[b]{\linewidth}\raggedright
σ
\end{minipage} & \begin{minipage}[b]{\linewidth}\raggedright
vs boruvka
\end{minipage} & \begin{minipage}[b]{\linewidth}\raggedright
vs bms
\end{minipage} \\
\midrule\noalign{}
\endhead
\bottomrule\noalign{}
\endlastfoot
10,000 & 50,000 & 25 & boruvka & 80.8 & 11.1 & --- & 0.41× \\
10,000 & 50,000 & 25 & bms & 33.1 & 14.3 & 2.44× & --- \\
10,000 & 50,000 & 25 & par\_T2 & 92.0 & 10.1 & 0.88× & 0.36× \\
10,000 & 50,000 & 25 & par\_T4 & 91.3 & 8.6 & 0.88× & 0.36× \\
100,000 & 500,000 & 6 & boruvka & 722.9 & 26.7 & --- & 0.37× \\
100,000 & 500,000 & 6 & bms & 268.3 & 13.1 & 2.69× & --- \\
100,000 & 500,000 & 6 & par\_T2 & 839.6 & 34.6 & 0.86× & 0.32× \\
100,000 & 500,000 & 6 & par\_T4 & 813.1 & 21.9 & 0.89× & 0.33× \\
500,000 & 2,500,000 & 3 & boruvka & 3870.6 & 78.7 & --- & 0.36× \\
500,000 & 2,500,000 & 3 & bms & 1404.3 & 2.0 & 2.76× & --- \\
500,000 & 2,500,000 & 3 & par\_T2 & 4109.9 & 16.2 & 0.94× & 0.34× \\
500,000 & 2,500,000 & 3 & par\_T4 & 4150.9 & 56.8 & 0.93× & 0.34× \\
1,000,000 & 5,000,000 & 3 & boruvka & 7950.6 & 61.5 & --- & 0.37× \\
1,000,000 & 5,000,000 & 3 & bms & 2921.4 & 62.6 & 2.72× & --- \\
1,000,000 & 5,000,000 & 3 & par\_T2 & 8556.6 & 47.6 & 0.93× & 0.34× \\
\end{longtable}

Table 6. Runtime results for dense graphs (m ≈ 50n).

\begin{longtable}[]{@{}
  >{\raggedright\arraybackslash}p{(\columnwidth - 14\tabcolsep) * \real{0.1395}}
  >{\raggedright\arraybackslash}p{(\columnwidth - 14\tabcolsep) * \real{0.1419}}
  >{\raggedright\arraybackslash}p{(\columnwidth - 14\tabcolsep) * \real{0.0841}}
  >{\raggedright\arraybackslash}p{(\columnwidth - 14\tabcolsep) * \real{0.1398}}
  >{\raggedright\arraybackslash}p{(\columnwidth - 14\tabcolsep) * \real{0.1449}}
  >{\raggedright\arraybackslash}p{(\columnwidth - 14\tabcolsep) * \real{0.0869}}
  >{\raggedright\arraybackslash}p{(\columnwidth - 14\tabcolsep) * \real{0.1571}}
  >{\raggedright\arraybackslash}p{(\columnwidth - 14\tabcolsep) * \real{0.1057}}@{}}
\toprule\noalign{}
\begin{minipage}[b]{\linewidth}\raggedright
n
\end{minipage} & \begin{minipage}[b]{\linewidth}\raggedright
m
\end{minipage} & \begin{minipage}[b]{\linewidth}\raggedright
trials
\end{minipage} & \begin{minipage}[b]{\linewidth}\raggedright
Algorithm
\end{minipage} & \begin{minipage}[b]{\linewidth}\raggedright
Mean (ms)
\end{minipage} & \begin{minipage}[b]{\linewidth}\raggedright
σ
\end{minipage} & \begin{minipage}[b]{\linewidth}\raggedright
vs boruvka
\end{minipage} & \begin{minipage}[b]{\linewidth}\raggedright
vs bms
\end{minipage} \\
\midrule\noalign{}
\endhead
\bottomrule\noalign{}
\endlastfoot
50,000 & 2,500,000 & 3 & boruvka & 3137.6 & 41.8 & --- & 0.37× \\
50,000 & 2,500,000 & 3 & bms & 1160.7 & 11.1 & 2.70× & --- \\
50,000 & 2,500,000 & 3 & par\_T2 & 3385.9 & 10.5 & 0.93× & 0.34× \\
50,000 & 2,500,000 & 3 & par\_T4 & 3423.9 & 50.7 & 0.92× & 0.34× \\
100,000 & 5,000,000 & 3 & boruvka & 6430.7 & 22.7 & --- & 0.39× \\
100,000 & 5,000,000 & 3 & bms & 2495.4 & 33.3 & 2.58× & --- \\
100,000 & 5,000,000 & 3 & par\_T2 & 6749.1 & 23.7 & 0.95× & 0.37× \\
\end{longtable}

\textbf{BMSBoruvka vs boruvka (sparse)}. BMSBoruvka achieves a speedup
of 2.44× to 2.76× over sequential Borůvka across all sparse instances.
The speedup is consistent and grows slightly with n, rising from 2.44 at
n = 10000 to 2.76 at n = 500000. This trend reflects the increasing
advantage of graph compression as the number of Borůvka rounds required
grows with n: while boruvka must scan all m edges in each of its five to
seven rounds, BMSBoruvka compresses the graph after each super-step and
subsequent rounds operate on a much smaller edge set.

\textbf{Parallel Borůvka vs sequential (sparse)}. The parallel variants
(par\_T2, par\_T4) are consistently slower than sequential Borůvka on
this two-core machine. Both par\_T2 and par\_T4 incur approximately
12--15\% overhead relative to sequential at n = 100000, shrinking to
6--7\% at n = 1000000. The overhead arises from three sources: (i)
allocation and initialisation of T × n thread-local best-edge arrays per
round, costing O(Tn) per super-step; (ii) the sequential Phase 2
reduction over all T × n entries; and (iii) OpenMP thread-launch
latency. With only two physical cores, the parallel workload in Phase 1
is insufficient to amortise these costs. On machines with 16 or more
cores, where Phase 1 can distribute m edge operations over many threads,
the parallel variant is expected to show positive speedup.

\textbf{Dense graphs}. On dense instances (m ≈ 50n), BMSBoruvka is still
2.58×2.70× faster than sequential Borůvka in absolute wall-clock time.
However, as discussed in Section 5.5, this speedup no longer corresponds
to a theoretical complexity advantage: on dense graphs Phase B degrades
to O(m log n) regardless of algorithm, so both boruvka and BMSBoruvka
are asymptotically equivalent. The practical speedup observed here stems
from BMSBoruvka's fewer total rounds (due to larger per-super-step
shrinkage) rather than from reduced edge-scan work per round.

\subsection{Throughput \(\widehat{\mathbf{t}}\) and Complexity Confirmation}

Table 7 reports \(\widehat{t}\) and \(\xi_{\mu}\) for the sparse series,
tracking both metrics as n grows.

Table 7. Shrinkage factor \(\xi_{\mu}\) and normalised throughput
\(\widehat{t}\) for sparse graphs.
\({\widehat{t}}_{bms} = \frac{n\log_{2}^{\frac{2}{3}}n}{T_{\text{mean}}}\);
\({\widehat{t}}_{boruvka} = \frac{m\log_{2}n}{T_{\text{mean}}}\).

\begin{longtable}[]{@{}
  >{\raggedright\arraybackslash}p{(\columnwidth - 10\tabcolsep) * \real{0.1496}}
  >{\raggedright\arraybackslash}p{(\columnwidth - 10\tabcolsep) * \real{0.1496}}
  >{\raggedright\arraybackslash}p{(\columnwidth - 10\tabcolsep) * \real{0.1442}}
  >{\raggedright\arraybackslash}p{(\columnwidth - 10\tabcolsep) * \real{0.1549}}
  >{\raggedright\arraybackslash}p{(\columnwidth - 10\tabcolsep) * \real{0.1923}}
  >{\raggedright\arraybackslash}p{(\columnwidth - 10\tabcolsep) * \real{0.2094}}@{}}
\toprule\noalign{}
\begin{minipage}[b]{\linewidth}\raggedright
n
\end{minipage} & \begin{minipage}[b]{\linewidth}\raggedright
\(\xi_{\mu}\)(bms)
\end{minipage} & \begin{minipage}[b]{\linewidth}\raggedright
\(\xi_{theory}\)(2ᵗ)
\end{minipage} & \begin{minipage}[b]{\linewidth}\raggedright
\(\xi_{\mu}\)(boruvka)
\end{minipage} & \begin{minipage}[b]{\linewidth}\raggedright
\({\widehat{t}}_{bms}\)(ops/ms)
\end{minipage} & \begin{minipage}[b]{\linewidth}\raggedright
\({\widehat{t}}_{boruvka}\)(ops/ms)
\end{minipage} \\
\midrule\noalign{}
\endhead
\bottomrule\noalign{}
\endlastfoot
10,000 & 40.82 & 64 & 5.07 & 1,695 & 8,223 \\
100,000 & 45.64 & 128 & 5.67 & 2,426 & 11,488 \\
500,000 & 44.19 & 256 & 5.35 & 2,529 & 12,228 \\
1,000,000 & 42.84 & 256 & 5.90 & 2,516 & 12,535 \\
\end{longtable}

\({\widehat{t}}_{boruvka}\) rises from 8,223 at n = 10000 to 12,535 at n
= 1000000, stabilising in the range 11,000--13,000 for n geq 100000.
This convergence confirms that the O(m log n) bound is tight: the
algorithm processes a consistent number of theoretical operations per
millisecond at scale, with no hidden super-linear growth term.

\({\widehat{t}}_{bms}\) rises from 1,695 to 2,516 and stabilises in the
range 2,400--2,530 for n geq 100000. This confirms that the
\(O\left( nlog^{\frac{2}{3}}n \right)\) bound is tight for BMSBoruvka on
sparse graphs. The lower absolute value of \({\widehat{t}}_{bms}\)
compared to \(\widehat{t}\)\_boruvka is not a contradiction: the two
normalisations use different theoretical operation counts
(\(nlog^{\frac{2}{3}}n\) versus m log n) that are not directly
comparable in magnitude.

The upward trend in both \(\widehat{t}\) values at small n is explained
by cache behaviour. At n = 10000, the DSU parent{[}{]} array occupies
only 40\,KIB and fits comfortably in L1/L2 cache; find() calls resolve
with minimal memory latency. As n grows beyond the cache boundary,
find() accesses become cache misses, increasing the real cost per
theoretical operation. By n = 100000 both effects have plateaued and
\(\widehat{t}\) stabilises, confirming that the RAM model captures the
asymptotic regime correctly even though it does not model cache effects.

Dense graph behaviour. For dense graphs, \({\widehat{t}}_{bms}\) drops
to approximately 260--270, far below the sparse-graph values. This is an
artefact of the normalisation: \({\widehat{t}}_{bms}\) uses
\(nlog^{\frac{2}{3}}n\) as its theoretical operation count, which is
derived under the sparse assumption m = O(n). On dense graphs Phase B
costs O(m log n) (Theorem 3), not
\(O\left( nlog^{\frac{2}{3}}n \right)\). Normalising by the smaller
sparse formula therefore makes \({\widehat{t}}_{bms}\) appear
artificially low. This is precisely the intended behaviour:
\({\widehat{t}}_{bms} \ll {\widehat{t}}_{boruvka}\) on dense graphs
serves as an empirical signal that the theoretical advantage of
BMSBoruvka does not apply, consistent with the dense-graph remark
following Theorem 3.

\subsection{Summary}

Table 8. Summary of experimental findings.

\begin{longtable}[]{@{}
  >{\raggedright\arraybackslash}p{(\columnwidth - 4\tabcolsep) * \real{0.3419}}
  >{\raggedright\arraybackslash}p{(\columnwidth - 4\tabcolsep) * \real{0.3312}}
  >{\raggedright\arraybackslash}p{(\columnwidth - 4\tabcolsep) * \real{0.3269}}@{}}
\toprule\noalign{}
\begin{minipage}[b]{\linewidth}\raggedright
Metric
\end{minipage} & \begin{minipage}[b]{\linewidth}\raggedright
Sparse graphs (m ≈ 5n)
\end{minipage} & \begin{minipage}[b]{\linewidth}\raggedright
Dense graphs (m ≈ 50n)
\end{minipage} \\
\midrule\noalign{}
\endhead
\bottomrule\noalign{}
\endlastfoot
BMSBoruvka speedup vs boruvka & 2.44× -- 2.76× & 2.58× -- 2.70× \\
Parallel Borůvka vs boruvka (2-core) & 0.86× -- 0.94× & 0.92× --
0.95× \\
\({\widehat{t}}_{bms}\) converges? & Yes, \textasciitilde2400--2530
(n≥1e5) & No (wrong normalisation basis) \\
\({\widehat{t}}_{boruvka}\) converges? & Yes,
\textasciitilde11000--13000 (n≥1e5) & Yes,
\textasciitilde12300--13000 \\
\(O(m\ log\ n)\) bound tight? & Confirmed (\({\widehat{t}}_{boruvka}\)
stable) & Confirmed (\({\widehat{t}}_{boruvka}\) stable) \\
\(O\left( nlog^{\frac{2}{3}}n \right)\) bound tight? & Confirmed
(\({\widehat{t}}_{bms}\) stable) & N/A (dense graph, m≠O(n)) \\
BMSBoruvka super-steps & 2 (all sizes) & 0--2 \\
boruvka rounds & 5--7 & 6 \\
\end{longtable}

The experiments confirm all major theoretical predictions. BMSBoruvka
achieves a consistent 2.5×--2.8× wall-clock speedup over sequential
Borůvka on sparse graphs, and the normalised throughput \(\widehat{t}\)
stabilises from n = 100000 onward for both algorithms, confirming that
\(O\left( nlog^{\frac{2}{3}}n \right)\) and \(O(m\ log\ n)\) are
empirically tight bounds. The parallel variant does not improve on
sequential performance at two cores due to per-round allocation overhead
and the serial reduction phase, but the underlying Phase 1 scan is
embarrassingly parallel and is expected to scale on machines with more
cores. On dense graphs, BMSBoruvka retains a practical speedup but loses
its theoretical complexity advantage, as predicted by the sparse-graph
condition \(m\  = \ O(n)\) established in Lemma 5.

\section{Discussion}

\subsection{Comparison with Classical MST Algorithms}

The minimum spanning tree problem has been studied for nearly a century,
and the evolution of its algorithms reflects a steady push towards lower
asymptotic complexity. Table~9 places BMSBoruvka in the context of the
major results in this history, across three dimensions: worst-case time
complexity, cost model, and practical implementability.

\hspace{0pt} Table 9. Selected MST algorithms in chronological order.
α = inverse Ackermann function; T*(m,n) = decision-tree complexity of
MST (exact value unknown).

\begin{longtable}[]{@{}
  >{\raggedright\arraybackslash}p{(\columnwidth - 8\tabcolsep) * \real{0.2998}}
  >{\raggedright\arraybackslash}p{(\columnwidth - 8\tabcolsep) * \real{0.2371}}
  >{\raggedright\arraybackslash}p{(\columnwidth - 8\tabcolsep) * \real{0.1593}}
  >{\raggedright\arraybackslash}p{(\columnwidth - 8\tabcolsep) * \real{0.1631}}
  >{\raggedright\arraybackslash}p{(\columnwidth - 8\tabcolsep) * \real{0.1407}}@{}}
\toprule\noalign{}
\begin{minipage}[b]{\linewidth}\raggedright
Algorithm
\end{minipage} & \begin{minipage}[b]{\linewidth}\raggedright
Time complexity
\end{minipage} & \begin{minipage}[b]{\linewidth}\raggedright
Model
\end{minipage} & \begin{minipage}[b]{\linewidth}\raggedright
Randomised?
\end{minipage} & \begin{minipage}[b]{\linewidth}\raggedright
Practical?
\end{minipage} \\
\midrule\noalign{}
\endhead
\bottomrule\noalign{}
\endlastfoot
Borůvka~{[}Bor26{]} & \(O(m\ log\ n)\) & Comparison & No & Yes \\
Kruskal~{[}Kru56{]} & \(O(m\ log\ n)\) & Comparison & No & Yes \\
Prim~{[}Pri57{]}~+~Fibonacci heap~{[}FT87{]} & \(O(m\  + \ n\ log\ n)\)
& Comparison & No & Yes \\
Karger, Klein \& Tarjan~{[}KKT95{]} & \(O(m)\) expected & Comparison &
Yes & Limited \\
Chazelle~{[}Cha00{]} & \(O(m\ \alpha(m,\ n))\) & Comparison & No & No \\
Pettie \& Ramachandran~{[}PR02{]} & \(O\left( T^{*}(m,n) \right)\) &
Comparison & No & No \\
BMSBoruvka (this work) & \(O\left( nlog^{\frac{2}{3}}n \right)\)
{[}sparse{]} & RAM & No & Yes \\
BMSBoruvka (this work) & \(O(m\ log\ n)\) {[}general{]} & RAM & No &
Yes \\
\end{longtable}

\textbf{Classical algorithms} (\textbf{Borůvka}~{[}Bor26{]},
\textbf{Kruskal}~{[}Kru56{]}, \textbf{Prim}~{[}Pri57{]}). The three
textbook algorithms all run in \(O(m\ log\ n)\) or \(O(m + n\ log\ n)\)
time under the comparison model. Borůvka's algorithm~{[}Bor26{]}, the
structural foundation of BMSBoruvka, iterates rounds in which every
component simultaneously selects its cheapest outgoing edge; after ⌈log₂
n⌉ rounds a single component remains. Kruskal's algorithm~{[}Kru56{]}
sorts all edges once and greedily adds them in order, with cost
dominated by the \(O(m\ log\ m)\) sort. Prim's algorithm~{[}Pri57{]},
when implemented with a Fibonacci heap~{[}FT87{]}, achieves the best
classical bound of \(O(m + n\ log\ n)\), which is tighter on dense
graphs but equivalent to Borůvka on sparse ones where \(m\  = \ O(n)\).
All three are straightforward to implement and have been standard in
pedagogy and practice for decades. BMSBoruvka improves on all three for
sparse graphs, achieving \(O\left( nlog^{\frac{2}{3}}n \right)\) versus
the best classical bound of \(O(n\ log\ n)\) when \(m\  = \ O(n)\), an
improvement by a factor of \({log}^{1/3}n\).

\textbf{Karger, Klein \& Tarjan} (1995)~{[}KKT95{]}. Karger, Klein, and
Tarjan gave the first linear-time MST algorithm, achieving \(O(m)\)
expected time under the comparison model by combining random edge
sampling with a linear-time MST verification subroutine. The algorithm
is optimal in expectation: no comparison-based algorithm can do better
than \(\Omega(m)\) in expectation. However, its reliance on randomness
is a meaningful drawback. The algorithm consumes \(O(m)\) random bits
and uses MST verification as a black box, making it substantially more
complex to implement correctly than the classical algorithms or
BMSBoruvka. In practice, the large constant factors and the overhead of
the verification step mean that the algorithm is rarely competitive on
realistic inputs despite its superior asymptotic bound. BMSBoruvka,
while asymptotically slower on general graphs, is deterministic,
requires no random bits, and is faster in practice on the sparse graph
inputs most common in engineering applications.

\textbf{Chazelle} (2000)~{[}Cha00{]}. Chazelle's algorithm is the best
known deterministic MST algorithm in the comparison model, running in
\(O(m\ \alpha(m,\ n))\) time, where α is the inverse Ackermann function.
Since \(\alpha(m,\ n) \leq 4\) for any input that fits in computer
memory, this is effectively \(O(m)\) in practice. The algorithm is built
around the soft heap, an approximate priority queue that deliberately
corrupts a small fraction of its keys in exchange for \(O(1)\) amortised
operation time. Despite its elegance, Chazelle's algorithm has never
been implemented in practice: the soft heap introduces enormous constant
factors, the algorithm is designed for a pointer machine rather than a
standard RAM, and the correctness argument is extremely delicate.
BMSBoruvka is asymptotically inferior to Chazelle on sparse graphs
(\(O\left( nlog^{\frac{2}{3}}n \right)\) vs \(O(m\ \alpha(m,\ n))\)) but
offers something Chazelle's algorithm does not: a working
implementation. Using only a DSU and std::nth\_element, BMSBoruvka
achieves \(O\left( nlog^{\frac{2}{3}}n \right)\) on sparse graphs with a
codebase under 400 lines of standard C++, verified correct on 3,000
random test cases.

\textbf{Pettie \& Ramachandran} (2002)~{[}PR02{]}. Pettie and
Ramachandran proved that the algorithmic complexity of MST equals its
decision-tree complexity, presenting a deterministic algorithm running
in \(O\left( T^{*}(m,n) \right)\) time, where \(T^{*}(m,n)\) is the
minimum number of edge-weight comparisons needed to determine the MST on
any graph with n vertices and m edges. This algorithm is provably
optimal in the comparison model. The current best bounds on \(T^{*}\)
are \(\Omega(m) \leq T^{*}(m,n) \leq O\left( m\alpha(m,n) \right)\), so
the exact complexity of MST remains unknown. The algorithm is purely
theoretical: it operates by precomputing optimal decision trees for all
graphs up to a fixed size r, and the constant factor is exponential in
r. No practical implementation exists or is expected. BMSBoruvka makes
no claim to theoretical optimality, but its
\(O\left( nlog^{\frac{2}{3}}n \right)\) sparse bound is derived
rigorously, and the algorithm runs correctly on inputs of up to
\(n = 10^{6}\) vertices in under 3 seconds on commodity hardware.

\subsection{Positioning of BMSBoruvka}

The comparison above reveals a clear division in the MST algorithm
landscape. On one side are the theoretically superior
algorithms---Karger-Klein-Tarjan~{[}KKT95{]} (\(O(m)\) expected) and
Chazelle~{[}Cha00{]} (\(O(m\ \alpha(m,\ n))\))---that achieve lower
asymptotic bounds but are either randomised, impractical to implement,
or both. On the other side are the classical practical
algorithms---Borůvka~{[}Bor26{]}, Kruskal~{[}Kru56{]}, and
Prim~{[}Pri57{]}---that are widely deployed but plateau at
\(O(m\ log\ n)\).

BMSBoruvka occupies a distinct middle position. For sparse graphs
(\(m\  = \ O(n)\)), it achieves \(O\left( nlog^{\frac{2}{3}}n \right)\),
which is asymptotically better than every classical algorithm and better
than Prim on sparse graphs by a factor of \(log^{\frac{1}{3}}n\). It is
deterministic, requires no random bits, and its implementation uses only
standard data structures available in any modern programming
environment. Empirically, it achieves 2.44× to 2.76× speedup over
sequential Borůvka on sparse inputs (Section~5), and the normalised
throughput \(\widehat{t}\) stabilises above n = 100,000, confirming that
the \(O\left( nlog^{\frac{2}{3}}n \right)\) bound is empirically tight.

On general graphs, BMSBoruvka degrades to \(O(m\ log\ n)\), matching
Borůvka~{[}Bor26{]}. This is expected: the \(log^{\frac{2}{3}}\)
advantage arises specifically from graph compression on sparse inputs,
where the compressed super-graph shrinks rapidly (Lemma~5). On dense
graphs, Phase B cannot reduce the edge set significantly through
compression, and the algorithm reverts to the same asymptotic regime as
its classical predecessors. This limitation is consistent with the
behaviour of all known sub-\(O(m\ log\ n)\) MST algorithms: they either
require randomness~({[}KKT95{]}) or are
impractical~({[}Cha00{]},~{[}PR02{]}) to transcend the
\(O(m\ log\ n)\ \)bound on dense inputs in a deployable form.

\section*{References}

\noindent \textbf{{[}Bor26{]}} Borůvka, O. (1926). O jistém problému minimálním.
\emph{Práce Moravské Přírodovědecké Společnosti}, 3, 37--58.
\smallskip


\noindent \textbf{{[}Cha00{]}} Chazelle, B. (2000). A minimum spanning tree
algorithm with inverse-Ackermann type complexity. \emph{Journal of the
ACM}, 47(6), 1028--1047. \url{https://doi.org/10.1145/355541.355562}
\smallskip


\noindent \textbf{{[}DMMS+25{]}} Duan, R., Mao, J., Mao, X., Shu, X., \& Yin, L.
(2025). \emph{Breaking the sorting barrier for directed single-source
shortest paths}. arXiv. \url{https://doi.org/10.48550/arXiv.2504.17033}
\smallskip


\noindent \textbf{{[}FT87{]}} Fredman, M. L., \& Tarjan, R. E. (1987). Fibonacci
heaps and their uses in improved network optimization algorithms.
\emph{Journal of the ACM}, 34(3), 596--615.
\url{https://doi.org/10.1145/28869.28874}
\smallskip


\noindent \textbf{{[}Hil24{]}} Hillyard, R. C. (2024). \emph{Data structures,
algorithms, and invariants: A practical guide}. Cognella Academic
Publishing. ISBN 978-1-7935-8884-5.
\smallskip


\noindent \textbf{{[}KKT95{]}} Karger, D. R., Klein, P. N., \& Tarjan, R. E.
(1995). A randomized linear-time algorithm to find minimum spanning
trees. \emph{Journal of the ACM}, 42(2), 321--328.
\url{https://doi.org/10.1145/201019.201022}
\smallskip


\noindent \textbf{{[}Kru56{]}} Kruskal, J. B. (1956). On the shortest spanning
subtree of a graph and the traveling salesman problem. \emph{Proceedings
of the American Mathematical Society}, 7(1), 48--50.
\url{https://doi.org/10.1090/S0002-9939-1956-0078686-7}
\smallskip


\noindent \textbf{{[}PR02{]}} Pettie, S., \& Ramachandran, V. (2002). An optimal
minimum spanning tree algorithm. \emph{Journal of the ACM}, 49(1),
16--34. \url{https://doi.org/10.1145/505241.505243}
\smallskip


\noindent \textbf{{[}Pri57{]}} Prim, R. C. (1957). Shortest connection networks
and some generalizations. \emph{Bell System Technical Journal}, 36(6),
1389--1401. \url{https://doi.org/10.1002/j.1538-7305.1957.tb01515.x}
\smallskip


\noindent \textbf{{[}Tar75{]}} Tarjan, R. E. (1975). Efficiency of a good but not
linear set union algorithm. \emph{Journal of the ACM}, 22(2), 215--225.
\url{https://doi.org/10.1145/321879.321884}
\smallskip


\noindent \textbf{{[}TvL84{]}} Tarjan, R. E., \& van Leeuwen, J. (1984).
Worst-case analysis of set union algorithms. \emph{Journal of the ACM},
31(2), 245--281. \url{https://doi.org/10.1145/62.2160}
\smallskip



\end{document}
